\scanpage{48}
\mspage{7}
% L1–5
Base ordonnée $(\alpha_{\unsure{i}})_{1\le i\le6}$ d'un graphe cubique : déduite de $(\unsure{\zeta|\xi}_i)_{1\le i\le6}$ dans $\Delta_0$ standard par \deleted{\unsure{la}} un isom. Base $(\alpha_i)_{i\in\underline I}$ : $\operatorname{Card}\underline I=6$, et on peut y ordonner \inserted{$\underline I,\ \underline I\simeq[1,6]$} de façon que la famille devienne une base ordonnée. Partie basique de $S$.

% L6–14
\uline{Prop} (\unsure{trivial|triviale}) $u\mapsto u((\zeta_i))$ définit une \deleted{isom} bijection de $\operatorname{Isom}(\Delta_0,\Delta)$ avec l'ens des bases ordonnées de $\Delta$ (qui est donc un torseur à droite sous $W_0=\operatorname{Aut}(\Delta_0)$, à gauche sous $W=\operatorname{Aut}(S)$). Considérons l'inclusion évidente $\unsure{\mathfrak S|\mathfrak G}_6\subset W_0$. Alors $u\mapsto u(\{\zeta_i\}_{i\in\underline I}$ définit par passage au quotient une bijection
\[
\operatorname{Isom}(\Delta_0,\Delta)/\mathfrak S_6\simeq
\begin{aligned}&\text{ens des parties}\\&\text{basiques de }\Delta.\end{aligned}
\]
% L15–24
\uline{Dém} Pour le premier pt, il suffit de voir que tt autom $u$ de $S$ tel que $u(\zeta_i)=\zeta_i$ \inserted{$\forall i$} est l'identité. Or comme $\zeta'_i$ est l'unique élément de $\Delta$ \deleted{distinct de $\zeta_i$ et des autres $\zeta_j$} qui \deleted{est} est lié à \deleted{aucun} \inserted{tous} des autres $\zeta_j$ \deleted{\unsure{t} \uline{NB} Cet élé} $u$ laisse invariant aussi les $\zeta'_i$, donc les $\zeta_{ij}$ ($\zeta_{ij}$ est l'unique élément de $S$ qui est lié à la fois à $\zeta_i$ et $\zeta'_j$, ou à la fois à $\zeta_j$ et $\zeta'_i$), donc est l'identité. \deleted{\unsure{Le|Ce}} Le $=$ argument montre que tt autom. de $\Delta_0$ qui \unsure{inverse|invarie|invar.} la partie $\{\zeta_i\}$ est induit par un $\sigma\in\mathfrak S_6$.
% archive sheet number: 47
